\documentclass[11pt,a4paper]{article}
\usepackage[T1]{fontenc}
\usepackage{lmodern,amsmath,amssymb,amsthm,mathtools,booktabs,microtype,xcolor}
\usepackage[margin=26mm]{geometry}
\usepackage[colorlinks=true,linkcolor=blue!45!black,citecolor=blue!45!black,urlcolor=blue!45!black]{hyperref}
\newtheorem{theorem}{Theorem}[section]\newtheorem{lemma}[theorem]{Lemma}
\newtheorem{proposition}[theorem]{Proposition}
\newcommand{\Q}{\mathbb Q}\newcommand{\Z}{\mathbb Z}\newcommand{\HH}{\mathcal H}
\DeclareMathOperator{\Tr}{Tr}\DeclareMathOperator{\Nm}{Nm}
\setlength{\emergencystretch}{2em}
\title{A family of defining quartics in 263 variables}
\author{}\date{9 November 2026}
\begin{document}\maketitle
\begin{abstract}
Under the stated arithmetic premises of the elliptic source, we specify
an integer quartic family in 263 total variables whose sufficiently
large members define the integers in the rationals. Four disjoint
affine images of the contact set give a height exponent $19/4$. This permits
quadratic Taylor approximations to the formal logarithm, with comparison
orders twelve and three. Residue classes modulo eight permit sixteen local tests to be represented
by five factors. A contact set symmetric about one makes its degree-twenty product
a degree-ten polynomial in the squared displacement. Five quadratic
coordinates evaluate that product and its normalization. The normalization
is independent of the scale and uses the positive aggregate trace directly. The complete
quadratic and expanded coefficient tables are independently verified.
The family has 3312 monomials for every positive parameter, and exact
coefficient height $80989^{4N+2}+1$ for $N\ge23$. The sufficient defining
parameter remains ineffective because the aggregate-trace height theorem
has an ineffective additive constant.
\end{abstract}

\section{The family and its arithmetic scope}
All variable counts include the distinguished input $X$. For a fixed
integer $N\ge1$ let $A_N=80989^{2N+1}$. The complete tables
\nolinkurl{quadratic-family263.json} and
\nolinkurl{polynomial-family263.json} specify 237 primitive integer
quadratic polynomials and their sum of squares
\[
 P_N(X,Y_1,\ldots,Y_{262})=\sum_{j=1}^{237}Q_{N,j}^2.
\]
The parameter $A_N$ is a coefficient, not a variable.

\begin{theorem}\label{thm:main}
Assume the arithmetic premises stated below. There is a fixed integer
$N_0\ge27464$ such that every $N\ge N_0$ satisfies
\[
 x\in\Z\quad\Longleftrightarrow\quad
 \exists\mathbf y\in\Q^{262}\ P_N(x,\mathbf y)=0
 \qquad(x\in\Q).
\]
For every $N\ge1$ the polynomial $P_N$ has degree exactly four,
content one, and 3312 nonzero coordinate monomials. For $N\ge23$,
\[
                       H(P_N)=80989^{4N+2}+1.
\]
The construction specifies every member exactly. The proof does not
determine a sufficient numerical value of $N_0$.
\end{theorem}

Set
\[
 K=\Q(\rho),\quad \rho^2=2,\quad D=(3-\rho)^2=11-6\rho,
 \qquad E:y^2=x^3-Dx,
\]
\[
 P=2(1+\rho,2\rho)
 =\left(\frac{57-28\rho}{8},\frac{-660+417\rho}{32}\right).
\]
We retain the following arithmetic premises, with the precise domains
and hypotheses recorded in \cite{Dependencies,Source,Weights}:
\begin{enumerate}
\item $2E(K)=\Z P$ and $h(x([n]P))\ge n^2-6$ for every nonzero
integer $n$, together with the specified point charts and group domains.
\item The five-contact height estimate with exponent nineteen,
exceptional cutoff $76176$, and the fixed constant
\[
 H=2^{110}210^{1800}\left(\prod_{p\le76176}p\right)^{19}+128.
\]
\item The pole obstruction outside $2,3,17$, and its simultaneous
realization for positive inputs less than $17$ with square reduced
denominator and the stated local squareclasses and valuation bounds
at those three primes. Its normalized quadratic equations are retained.
\item The simultaneous prime-pattern, conjugate reduction-order,
finite congruence, and real approximation assertions used by the source.
The weighted selection and fixed denominator dilation are justified
from the cited analytic estimates in \cite{Selection}.
\end{enumerate}
The aggregate trace-height theorem used below is the proved consequence
of the family height theorem described in \cite{Aggregate}; its additive
constant is ineffective. The contact amplification in \cite{Contacts}
is an elementary consequence of the retained five-contact estimate,
not an additional arithmetic assumption. Exact coefficient verification
does not prove the retained arithmetic premises.

\section{Four contact blocks and the two scales}
Put
\[
 \Theta=\{0,-1/2,-1,1/2399,-32/407\},\qquad
 \Theta_4=\Theta\cup(2-\Theta)\cup(\Theta+5)\cup(-3-\Theta),
\]
\[
                         F_4(t)=\prod_{\theta\in\Theta_4}(t-\theta).
\]
The twenty contacts are distinct, omit one, and are symmetric about one.
Let $\Sigma$ consist of the primes at most $76176$ and the three primes
$247957,368323,2029147$. Outside this set, all contact differences and
all numbers $1-\theta$ are units. Define $M_4(t)$ as the product of primes
$p\notin\Sigma$ for which $v_p(t-\theta)$ is positive and odd for some
$\theta\in\Theta_4$. Exact factorizations in the independent certificate
verify the exceptional set. They also show $M_4(1)=1$.

Put $E_4=247957\cdot368323\cdot2029147$. Then
\begin{equation}\label{eq:contactheight}
 h(t)\le\log14+HE_4^{19}M_4(t)^{19/4}
       \le H_4M_4(t)^{19/4},\qquad
 H_4=\log14+HE_4^{19},
\end{equation}
for $t\notin\Theta_4$. Here is a direct proof from the retained
five-contact inequality. Apply it to the four maps
\[
 2+\frac1t,\qquad 2+\frac1{2-t},\qquad
 2+\frac1{t-5},\qquad2+\frac1{-3-t}.
\]
Their integer matrices have determinant $\pm1$. The maximum absolute
row sums of inverse projective matrices are respectively $3,7,14,8$,
so each inverse changes logarithmic height by at most $\log14$.
Let $M_j$ be the contact radical for the corresponding five-element
block, omitting $\Sigma$. Its original radical is at most $E_4M_j$.
Thus $h(t)-\log14\le HE_4^{19}M_j^{19}$ for each of the four blocks.
If the left side is positive, multiplication gives the first inequality
in~\eqref{eq:contactheight}, since the block radicals are pairwise
coprime and their product is $M_4(t)$. If it is nonpositive, the
inequality is immediate. This accounts for every fixed exceptional prime.

The retained root points are
\[
 Q_1=[n]P,\quad Q_2=[a]P,\quad Q_3=[t]P,\quad Q_4=[b]P,
 \qquad B=\max(|n|,|a|,|t|,|b|),
\]
with nonzero integer labels. Set
\[
 T=\frac14\Tr_{K/\Q}\left(\sum_{i=1}^4x(Q_i)\right).
\]
It is positive, and the aggregate trace theorem gives fixed constants
$M_1>0$, $M_2\ge1$ such that
\begin{equation}\label{eq:aggregate}
                         B^2\le M_1h(T)+M_2.
\end{equation}
Use the two positive rational scales $a_*=1,T$. No square condition
on a scale is needed. Replacing the half-trace in the aggregate
height theorem by its half changes its logarithmic height by at
most $\log2$; this is absorbed into $M_2$.
For each scale put $t_*=a_*\lambda$ and $\delta=t_*-1$.
On the nonzero domain the detector is
\begin{equation}\label{eq:detector}
 R=\frac{F_4(t_*)}{\delta^{22}},\qquad
 q=\frac{\delta^{18}}{F_4(t_*)}=\frac1{R\delta^4}.
\end{equation}
The contact circuit stores $r=\delta^2$ and its powers $r^2,r^4,r^5$.
The actual quadratic normalization equations are
\begin{equation}\label{eq:clearing}
                         F_4(t_*)q=r^4r^5,\qquad k=qr^2.
\end{equation}
Thus $k=R^{-1}$ when $\delta\ne0$. No contact value can occur:
$F_4(t_*)=0$ would force $r=0$ and $t_*=1$, contradicting $F_4(1)\ne0$.
If $q=0$, the same equation forces $\delta=0$. At $\delta=0$ the
contact radical is one; no reciprocal is used on this stratum.

If both radicals divide a nonzero integer $J$ to order $k$, then
$h(T)\le h(\lambda)+h(T\lambda)$ and
\eqref{eq:contactheight}--\eqref{eq:aggregate} imply
\begin{equation}\label{eq:heighttransfer}
 B^2\le C|J|^{19/(4k)}+M_2,\qquad C=2M_1H_4.
\end{equation}
The two radicals need not be coprime for this implication.

\section{Local tests and the approximation denominators}
Let $\sigma$ be conjugation, $f=1/x$, and $z=-x/y$, extended by zero
at the identity where these functions are regular. The retained
point relations are
\[
 U_0=H,\quad V_0=H-P,\quad U_i=[u_i]P,\quad V_i=U_i-Q_i.
\]
The charts permit $U_i=O$ or $V_i=O$ for $i>0$.
Put
\[
 G_2(w)=1-\frac25w-\frac2{15}w^2,
 \quad \ell_0=z(V_0)G_2(Df(V_0)^2),\quad
 L_i=\frac{z(U_i)}{z(H)},
\]
\[
 L_1'=\sigma\!\left(\frac{z(V_1)}{\ell_0}\right),\qquad
 L_i'=\sigma\!\left(\frac{z(V_i)G_2(Df(V_i)^2)}{\ell_0}\right)\quad(i>1).
\]
These quotients are interpreted only after the denominator exclusions
below. Write
\[
 D_H=f(H),\quad A_i=f(U_i),\quad B_i=\sigma f(V_i),\quad
 C_i=L_i-L_i',\quad I=L_1-\beta,\quad \beta=51X^2+2504,
\]
\[
 G=L_3^2-L_2(L_2-2)L_3L_4+(L_2-2)^2L_4^2-1.
\]
The required orders are six for $D_H$, twelve for $C_2,C_3,C_4,G$,
three for $C_1$, six for $I$, and one for the nine $A_i,B_i$.
The final input comparison uses order three, the minimum of the two
orders controlling $L_1-n$ and $L_1-\beta$.

\begin{lemma}\label{lem:g2}
The polynomial $G_2$ has no zero in $K$.
\end{lemma}
\begin{proof}
Its roots are $(-3\pm\sqrt{39})/2$, and $39$ is neither a rational
square nor twice a rational square. Explicitly, if
$(u+v\sqrt2)^2=39$ with $u,v\in\Q$, then $uv=0$; the two cases
would require $u^2=39$ or $v^2=39/2$, both impossible.
The independent finite certificate also verifies that the integer
numerator $15-6w-2w^2$ has no root modulo 17, a prime that splits in
$K$. Its leading coefficient is a unit there, so a hypothetical local
root would be integral and give a forbidden residue root.
\end{proof}

The finite $H$ chart is $xf=1$, $S^2=x-Df$, $H=(x,xS)$.
The source group calculation writes $H=[m]R_0+T_0$ with $P=2R_0$,
$m\in\Z$, and $T_0\in E[2](K)$. Put $h=m/2$ for its rational
free index. The exceptional classes $E[2](K),\pm P+E[2](K)$ are
excluded before any logarithmic ratio is interpreted. Indeed their
finite allowed reciprocal abscissas are units outside the fixed small
primes, whereas extraction below would require positive valuation
at any contact. All radicals would then be one, forcing
$B^2\le C+M_2$ by \eqref{eq:heighttransfer}. The retained scalar
lower bound in Section~\ref{sec:soundness} contradicts this for every
sufficiently large parameter. The explicit torsion translations and
norm values supporting the unit assertion are those recorded in
\cite[Section 3]{Even}. This exclusion uses no comparison quotient.

Consequently $m\notin\{-2,0,2\}$, and $H,V_0$ are nontorsion. Their required
nonzero abscissas and ordinates, and all subtraction denominators,
follow from the retained chart identities. In particular
$z(H),z(V_0)\ne0$. Lemma~\ref{lem:g2} gives $\ell_0\ne0$ globally.
The analytic formal logarithms have a separate nonvanishing proof
from the formal logarithm isometry when the points are formal.

\section{Five factors and extraction at one place}
At a contact prime $p\notin\Sigma$, $\delta$ is a unit and
$\nu=v_p(R)=v_p(F_4(t_*))$ is positive and odd. This follows
from the unique contact factor with positive odd valuation and
does not require any condition on the valuation of the scale. All valuations
of $K$ are unramified and normalized by $v(p)=1$.
For a point on $E$, every finite value of $v(f)$ is even: if
$v(x)<0$ then $2v(y)=3v(x)$, while if $v(x)>0$ then
$2v(y)=v(x)$; the remaining case is valuation zero.
This also holds for conjugated reciprocal abscissas. Before excluding
the exceptional $H$ classes, the stored value for $f(H-P)$ still has even
valuation: its literal equation is $f(H)F_0=g^2$, and $f(H)$ has even
valuation. Thus the parity argument and the preceding exclusion do not
assume the validity of a division in the fixed subtraction chart.

Use
\begin{align*}
 T_1&=1+q^{23}C_2^2+q^{10}D_H^2+q^4B_1^4,\\
 T_2&=1+q^{23}G^2+q^2A_3^2+q^4B_0^4,\\
 T_3&=1+q^{23}I^4+qA_2^2+q^2A_4^2+q^4B_3^4,\\
 T_4&=1+q^{23}C_3^2+q^{10}C_1^4+q^4B_4^4,\\
 T_5&=1+q^{23}C_4^2+q^2A_1^2+q^4B_2^4,
 \qquad W=T_1T_2T_3T_4T_5.
\end{align*}

\begin{lemma}\label{lem:residues}
At a contact prime, every factor is nonzero and has valuation at most
zero. A factor of valuation zero is congruent to one in the residue field.
\end{lemma}
\begin{proof}
First suppose every displayed test in a factor is nonzero. In each
factor other than $T_3$, the three nonconstant valuations belong,
respectively, to the odd integers, to $2+4\Z$, and to $4+8\Z$.
For example $T_1$ has valuations
\[
 2v(C_2)-23\nu,\qquad 2v(D_H)-10\nu,\qquad4v(B_1)-4\nu.
\]
The second has exact factor two because $v(D_H)$ is even and $\nu$ is
odd. The third has exact factor four because $v(B_1)$ is even.
The term $q^{10}C_1^4$ in $T_4$ has valuation
$2(2v(C_1)-5\nu)$, again twice an odd integer.

In $T_3$ the four nonconstant valuations are
\[
 4v(I)-23\nu,\quad 2v(A_2)-\nu,\quad
 2v(A_4)-2\nu,\quad4v(B_3)-4\nu.
\]
The first three have residues $\nu,-\nu,2$ modulo four. These are
pairwise distinct and nonzero because $\nu$ is odd and the $A_i$
have even valuations. The fourth has residue four modulo eight;
it therefore differs from each of the first three and from zero.
Thus, in every factor, all finite term valuations, including zero
for the constant term, are distinct. Zero tests simply remove terms
from this argument. The minimum is attained exactly once. This proves
nonvanishing and the asserted bound. If the minimum is zero, every
nonconstant term has positive valuation, so the residue is one.
\end{proof}

\begin{proposition}\label{prop:extract}
At each contact prime, the retained pole system and the cleared trace
equations imply $\nu>0$ and give all sixteen required test orders at
one place above that prime.
\end{proposition}
\begin{proof}
Write $W=w_0+w_1\rho$. The literal cleared trace equations are
\[
 w_0^2-2w_0\tau-2w_1^2=0,\qquad k\mathfrak b=\tau^2,\qquad k=R^{-1}.
\]
By Lemma~\ref{lem:residues}, $W\ne0$. The first equation therefore
gives $w_0\ne0$ and $\tau=1/\Tr(W^{-1})\ne0$; indeed $w_0=0$
would give $w_1=0$. Hence
$\mathfrak b=R/\Tr(W^{-1})^2$. If
$\ell=\max_{v\mid p}v(W)\le0$, the trace inequality gives
$v_p(\mathfrak b)\le\nu+2\ell$.
This valuation is odd; the retained pole obstruction forces it positive.
Consequently $\nu>0$ and $\ell>-\nu/2$.
At a place attaining $\ell$, the factor valuations are nonpositive
and add to $\ell$. Each factor valuation is therefore at least $\ell$,
and each of its term valuations is strictly greater than $-\nu/2$.

The resulting lower bounds, before integer rounding, are
\[
\begin{array}{c|c|c}
\text{test}&\text{strict lower bound}&\text{integer order}\\\hline
C_2,C_3,C_4,G&45\nu/4&12\\
I&45\nu/8&6\\
C_1&19\nu/8&3\\
D_H&19\nu/4&6\\
A_2&\nu/4&1\\
A_1,A_3,A_4&3\nu/4&1\\
B_0,B_1,B_2,B_3,B_4&7\nu/8&1
\end{array}
\]
Here $\nu$ is a positive odd integer. The order six for $D_H$ uses
its even valuation. In fact each of the nine reciprocal-abscissa
tests has order at least two; order one suffices for its subsequent
formal-group use. Zero tests satisfy all the stated inequalities.
\end{proof}

\subsection{Integrality precedes the depth estimate}
The extracted positive $f$ valuations put $H,U_i$ on the primary
formal branch and $V_0,V_i$ on the conjugate formal branch.
At a nonidentity formal point of depth $r=v(z)\ge1$, write
\[
 \log_E(Q)=z(Q)\mathcal A(Q),\qquad
 \mathcal A(Q)=G(Df(Q)^2),\qquad \mathcal A(O)=1.
\]
The formal series satisfies
\[
 G(w)=1-\frac25w-\frac2{15}w^2+O(w^3),\qquad
 g_j=\frac{4j-5}{4j+1}g_{j-1}\quad(j\ge2).
\]
The coefficient bound and full-tail proof in \cite[Lemma 3.1]{Even}
give $\mathcal A\equiv1\pmod{p^{4r}}$,
$G_2/\mathcal A\equiv1\pmod{p^{12r}}$, and
$1/\mathcal A\equiv1\pmod{p^{4r}}$.
These analytic units equal one at $O$, so zero numerator points
require no quotient $\log_E(O)/z(O)$.

The integral group identities
\[
 [m]U_i=[2u_i]H,\qquad
 [m-2]V_i=[2(u_i-n_i)]V_0
\]
kill the torsion component of $H$ and yield full logarithm ratios
$r_i=u_i/h$ and $r_i'=(u_i-n_i)/(h-1)$.
At the odd prime, $h,h-1$ are integral and at least one is a unit.
Hence at least one full ratio is integral. Multiplication by the
analytic units preserves the corresponding approximate valuation.
The positive-order cross test transfers integrality to the other
approximate ratio and thus to both full ratios.

Only now use $L_i=z(U_i)/z(H)$ integral. The $H$ test gives
$v(z(H))\ge3$, so every nonzero $U_i$ has depth at least three.
The primary errors have order at least twelve. The alternate errors
have order at least four for $i=1$ and twelve for $i>1$.
The integral full ratios convert these relative bounds into absolute
congruences. Multiplication by $h,h-1$ and subtraction then give
\[
 L_1\equiv n\pmod{p^3},\qquad
 L_i\equiv n_i\pmod{p^{12}}\quad(i>1).
\]
This argument includes zero $U_i$ or $V_i$, half-integral $h$, and
inert primes. At the inert primes conjugation preserves the valuation.

\section{Global soundness for sufficiently large parameters}\label{sec:soundness}
The retained scalar circuit gives, for $c=|a-1|$,
\begin{equation}\label{eq:scalar}
 c\ge80989^{L_N},\qquad
 L_N=\left\lfloor\frac{187(2N-1)}8\right\rfloor+1,
 \qquad h(\beta)\le\HH c^{19/24}.
\end{equation}
These statements, including their integral group identities and
exceptional-point exclusions, are the scalar conclusions of
\cite{Source}; no comparison logarithm is used to establish them.
Their circuit is retained, including the actual positive contact
product $r$ and the normalizer
\[
 r=\beta(\beta-1)(\beta-2401)(32\beta+343)(2\beta-1),
 \qquad R_s=\frac{A_Nr}{1+80989^2r}.
\]

Let
$\Delta=t^2-a(a-2)tb+(a-2)^2b^2-1$.
The local comparison shows that each twelfth power of a comparison
radical divides $\Delta$. If $\Delta\ne0$, then
$|\Delta|\le14B^4$, and \eqref{eq:heighttransfer} gives
\[
                  B^2\le C14^{19/48}B^{19/12}+M_2.
\]
Since $19/12<2$, this bounds $B$ by a fixed constant. Equation
\eqref{eq:scalar} gives $B\ge c-1\to\infty$ uniformly as
$N\to\infty$, a contradiction. Thus $\Delta=0$ eventually.
The retained integral Lucas descent then gives
\[
                         \log B>|a|-3\ge c-4.
\]

If $n-\beta\ne0$, let $J$ be its nonzero reduced numerator.
The local input comparisons make both cubed radicals divide $J$,
while $|J|\le2B\exp(h(\beta))$. For $B^2\ge2M_2$,
\eqref{eq:heighttransfer} gives
\[
 \log B\le\frac{12}{5}\log(2C)+\frac{19}{5}\log2
                +\frac{19}{5}h(\beta)
 \le C'+\frac{19}{5}\HH c^{19/24}.
\]
The right side is $o(c)$, contradicting the preceding lower bound.
Hence $n=\beta\in\Z$. If $X=u/v$ is reduced with $v>0$, then
$v^2\mid51$, and squarefreeness of $51$ gives $v=1$.
All bounds depend only on fixed constants. Since $L_N$ is increasing,
a single $N_0$ works for every larger $N$. Its dependence on the
aggregate additive constant is ineffective.

\section{Simultaneous completeness}
Fix an integer input and the coefficient parameter. Retain the scalar
and Lucas construction of the source, and choose the integer labels,
root points, two scales, and compact real and finite-place neighborhoods
before selecting primes. The neighborhoods avoid every retained
rational denominator; Lemma~\ref{lem:g2} ensures that the new
denominator introduces none. Every local test then has a fixed lower
valuation bound at each prime in the finite preparation set.

Choose a fixed positive integer $D_0$ clearing all contact coefficients
$D_0^2\theta/a_*$ and sufficiently divisible at the prepared primes,
including $\Sigma$, and the prime divisors of the scales. In the prime
pattern write
\[
                         \lambda=\frac{u}{D_0^2w^2}.
\]
The zero contact makes $u$ one of the selected primes, outside that
finite set. Thus $v_p(\lambda)\le-2v_p(D_0)$ there. Make this upper
bound low enough that $t_*=a_*\lambda$ has valuation below zero and
all nonzero contacts. Then $v_p(\delta)=v_p(t_*)$, and the exact
valuation formulas are
\[
 v_p(F_4(t_*))=20v_p(t_*),\qquad
 v_p(q)=v_p(R)=-2v_p(t_*).
\]
The last two valuations become uniformly arbitrarily large. Thus
all weighted test terms are small at every prepared place, and
$\Tr(W^{-1})$ is arbitrarily close to two.

The local squareclasses follow directly from
\begin{equation}\label{eq:preparedsquare}
 R=\delta^{-2}\prod_{\theta\in\Theta_4}
                 \left(1+\frac{1-\theta}{\delta}\right).
\end{equation}
The factors in the product are arbitrarily close to one. Requiring
congruence to one modulo eight at $2$, and modulo $p$ at an odd
prepared prime $p$, makes their product a local square. The prefactor
$\delta^{-2}$ is a square, as is the reciprocal trace factor in the
pole input. Taking $-v_p(t_*)$ sufficiently large also gives the
required pole valuations at $2,3,17$, including the subtraction
of $2v_p(2)$ in the trace denominator.

This preparation is valid for any fixed nonzero rational scales.
The prime-pattern input in \cite{Selection,Dependencies} applies to
arbitrary fixed distinct integer coefficients. After clearing
$D_0^2\theta/a_*$ and removing repeated coefficients, it applies
here without a square restriction on $a_*$. All denominator
conditions are fixed before prime selection.

Select the enlarged finite set of contact primes with the inherited
opposite-order separation. The common-index construction takes $H=[h]P$,
$U_i=[hn_i]P$, and $V_i=[(h-1)n_i]P$. Its congruences retain primary
depth twenty for $H$ and every nonzero $U_i$. Thus the full primary
logarithm ratio is exactly $n_i$, and its parameter-ratio approximation
has error of order eighty. In particular, since $n_1=\beta$ in the
completeness construction, $I=L_1-\beta$ has order eighty; the stronger
requirement six on this test imposes no new arithmetic condition. At the opposite places the points are formal. The new
Taylor approximations give comparison precision twelve, and the
constant approximation at $V_1$ gives precision four. The reciprocal-abscissa tests have valuation at least two. At the
chosen primary place, where $v_p(R)=1$, it therefore suffices to use
orders $6,12,4,12$ respectively for $D_H$, the high comparisons,
$C_1$, and $I$: every nonconstant term in all five factors has
positive valuation. All these orders follow from the preceding
construction. There $W\equiv1$. At the opposite
place, the valuation-class argument gives either positive valuation
for $W^{-1}$ or residue one. Its trace is consequently a unit with
residue one or two, including the inert case. Hence the pole input
has valuation one at every selected contact prime.

The exact prime-pattern identity is
\[
 R=\frac{\displaystyle a_*^{20}\prod_{\theta\in\Theta_4}
          (u-(\theta/a_*)D_0^2w^2)}
        {(D_0^2w^2)^{20}\delta^{22}}.
\]
Every factor outside the selected prime numerators is a rational
square, including $a_*^{20}$ even when $a_*$ is not itself a square.
Thus at every prime outside those numerators the valuation is even. Division by the squared trace
preserves this parity. This argument includes primes dividing $\delta$;
no unit assertion about $\delta$ is made away from contacts.
Equation~\eqref{eq:preparedsquare} and the preparatory divisibility
give the required local squareclasses and valuation bounds at $2,3,17$.

On the fixed compact real neighborhoods,
\[
 R\sim a_*^{-2}\lambda^{-2},\qquad
 q\sim a_*^{-2}\lambda^{-2},\qquad
 W\longrightarrow1,\qquad\Tr(W^{-1})\longrightarrow2.
\]
All test powers in the factors are even, so the pole input tends to
zero through positive values. It eventually lies in $(0,17)$.
The simultaneous selection supplies the required real size after all
fixed data have been chosen. The pole-realization input supplies its
witnesses, and the exact power identities fill the quadratic circuit.
This proves completeness with no additional condition on $N$.

\section{Exact quadratic specification}
The weights are specified by six scalar quadratic equations:
\[
 q_2=q^2,\quad q_4=q_2^2,\quad c q_4=q,\quad
 q_5=qq_4,\quad q_{10}=q_5^2,\quad c q_{23}=q_{10}^2.
\]
At every contact $q\ne0$, so $c=q^{-3}$ and each $q_j$ is exactly
$q^j$. The equations make no division by a variable. If $q=0$, the
clearing equations force $\delta=0$ and the contact radical is one,
as already proved. No value of the unconstrained reciprocal weight
on that stratum affects the soundness argument.

\begin{lemma}\label{lem:contactcircuit}
The values $F_4(t_*)$, $(t_*-1)^{18}$ and $(t_*-1)^4$ are quadratic
expressions in five additional coordinates, each determined by one
successive quadratic equation over $\Q$.
\end{lemma}
\begin{proof}
Put $r=(t_*-1)^2$. Symmetry of the contacts gives
\[
 F_4(t_*)=P_{10}(r),\qquad
 P_{10}(r)=\prod_{\theta\in\Theta}
          \bigl(r-(\theta-1)^2\bigr)\bigl(r-(\theta+4)^2\bigr).
\]
The ten squared roots are distinct. Four successive equations compute
\[
                  r=(t_*-1)^2,\quad r_2=r^2,\quad
                  r_4=r_2^2,\quad r_5=rr_4.
\]
Every integer from zero through ten is a sum of two elements of
$\{0,1,2,4,5\}$. Therefore $P_{10}(r)$ is a quadratic expression
in $1,r,r_2,r_4,r_5$, and a fifth equation computes it.
The other required values are already quadratic expressions:
\[
                    (t_*-1)^{18}=r_4r_5,\qquad
                    (t_*-1)^4=r_2.
\]
All coefficients are fixed rational numbers. These equations determine
the claimed values at every input, including the contact values and one.
\end{proof}

The normalization is consequently $F_4(t_*)q=r_4r_5$, followed by
$k=qr_2$. The trace-clearing coordinate is $k=R^{-1}$ on the nonzero
domain. The second scale is an affine rational expression in the retained
root-abscissa coordinates, so its product with $\lambda$ is quadratic
without a separate scale coordinate.

The high-point approximants $G_2$ are affine expressions in the
retained $f^2$ and $w^2$ coordinates. They introduce no additional
field value. At $V_1$ the numerator is simply $z(V_1)$.
The shared field powers requiring additional coordinates are
\[
                              C_1^4,\quad I^4,\quad B_1^4,
\]
each with two rational coordinates. The fourth powers of
$B_0,B_2,B_3,B_4$ are affine expressions in the Taylor circuits.
Each detector has five field factors and three private intermediate
products, using sixteen coordinates; the final product uses the
retained $W$ coordinates. The exact count is
\[
\begin{array}{lr}\toprule
\text{point, scalar, contact, logarithm, and weight coordinates}&225\\
\text{three shared fourth powers}&6\\
\text{five factors and intermediate products for two detectors}&32\\\midrule
\text{total coordinates, including the input}&263\\\midrule
\text{point, scalar, contact, logarithm, and weight quadratic rows}&195\\
\text{shared fourth-power rows}&6\\
\text{factor and product rows}&36\\\midrule
\text{total quadratic rows}&237\\\bottomrule
\end{array}
\]

The complete tables fix every coordinate, equation, coefficient, and
normalization; the input is coordinate one. The producer clears
denominators and primitive-normalizes each row before squaring.
An independent checker imports no producer. It recovers $P_{10}$ by
Newton finite differences at eleven exact arguments, checks all ten roots,
and independently verifies $P_{10}((t-1)^2)=F_4(t)$ as a degree-twenty
polynomial identity. It reconstructs every modified quadratic equation,
compares the other equations with their pinned coefficient source, and
checks the full contact exceptional set and all four height matrices.
Both normalizations and the linear trace scale are verified directly.
The projection proof uses these literal equations and the preceding
arithmetic argument.
It expands the entire sum of squares independently by unordered
pairs of terms. The supplementary local certificate exhausts all
odd residues of $\nu$ modulo eight and all possible valuation residues
of every test. These complete residue classes verify the
noncancellation assertion for all integer valuations, together with
the exact extraction bounds and sufficient completeness depths.
The infinite arithmetic arguments are the proofs above.

The expanded quartic has 3313 terms before powers of $A_N$ belonging
to the same coordinate monomial are collected, and 3312 monomials
after collection. Its unique coefficient involving $A_N^2$ is
$A_N^2+1$, at the square of the actual scalar contact-product coordinate.
Every other coefficient is a nonzero constant or a nonzero constant
multiple of $A_N$. Consequently no support cancellation occurs for
positive $A_N$. The largest constant coefficient $C_0$ has 1520 bits,
the largest linear multiplier is $13118436242$, and exact comparison
gives
\[
 C_0>13118436242^2,\qquad
 80989^{90}+1<C_0<80989^{94}+1.
\]
Thus $H(P_N)=\max(C_0,A_N^2+1)$ for every $N\ge1$ and the asserted
exact height holds for $N\ge23$. This is a coefficient statement;
it does not make the defining threshold effective.

\begin{thebibliography}{99}\small\setlength{\itemsep}{2pt}
\bibitem{Dependencies} Project repository,
\emph{Arithmetic dependencies of the defining-polynomial constructions},
\nolinkurl{arithmetic_dependencies_20261108_v8.tex}.
\bibitem{Source} Project repository,
\emph{A family of quartics in 297 variables},
\nolinkurl{quartic297_20261109_v12.tex}, especially Sections 2--6.
\bibitem{Weights} Project repository,
\emph{Common local weights and quartics in 297 and 342 variables},
\nolinkurl{shared_local_weights_20261109_v12.tex}.
\bibitem{Even} Project repository,
\emph{An explicit quartic in 322 variables},
\nolinkurl{quartic322_20261109_v13.tex}.
\bibitem{Contacts} Project repository,
\emph{Height Bounds from Disjoint Sets of Rational Contacts},
\nolinkurl{contact_height_amplification_20261109_v13.tex}.
\bibitem{Aggregate} Project repository,
\emph{The height of an aggregate elliptic trace},
\nolinkurl{aggregate_trace_height_20261108_v9.tex}.
\bibitem{Selection} Project repository,
\emph{Prime selection with fixed congruences and separated reduction orders},
\nolinkurl{prime_selection_20261109_v12.tex}.
\end{thebibliography}
\end{document}
